\rtmtight
\begin{tblr}{width=\textwidth,colspec={|Q[c,m,2.1cm]|X[l,m]|},hlines,vlines,hspan=minimal,rowsep=1.5pt,colsep=3pt,row{1}={bg=headgray,font=\bfseries}}
\SetCell{c,m}\hfil\logo\hfil & \textbf{POLITEKNIK NEGERI MALANG}\par\textbf{JURUSAN TEKNOLOGI INFORMASI}\par\textbf{PROGRAM STUDI: D4 TEKNIK INFORMATIKA} \\
\SetCell[c=2]{l} \textbf{RENCANA TUGAS MAHASISWA (RTM) DAN RUBRIK PENILAIAN} & \\
Nama Mata Kuliah & Kecerdasan Artifisial \\
Kode Mata Kuliah & RTI254004 \quad \textbf{sks / Jam perminggu:} 2 SKS/4 jam \quad \textbf{Semester:} 4 \\
Dosen Pengampu & Tim Pengajar Program Studi D4 TI \\
\end{tblr}
\assessmentblock{Tugas Implementasi Algoritma AI}{Case Method dan praktik pemrograman}{SCPMK0706-03101, SCPMK1005-03102}{Mahasiswa menyelesaikan tugas mengimplementasikan dan mengevaluasi algoritma AI sesuai Sub-CPMK dan konteks mata kuliah.}{Menganalisis kebutuhan, mengimplementasikan algoritma, menguji dan mengevaluasi kinerja, mendokumentasikan hasil.}{Kode sumber, laporan analisis, dan/atau bahan presentasi.}{Ketepatan implementasi algoritma & 12 & Algoritma diimplementasikan dengan benar sesuai spesifikasi, menghasilkan output yang tepat pada berbagai kasus uji. & Implementasi sebagian besar benar, tetapi ada kasus uji yang gagal atau efisiensi belum optimal. & Implementasi tidak berjalan dengan benar atau tidak sesuai algoritma yang diminta. \\ \\
Kualitas evaluasi dan analisis kinerja & 10 & Kinerja algoritma dievaluasi menggunakan metrik yang tepat, analisis mendalam, dan perbandingan antaralgoritma dilakukan secara valid. & Evaluasi dilakukan, tetapi metrik atau analisis perbandingan masih terbatas. & Tidak ada evaluasi kinerja yang bermakna atau metrik yang digunakan tidak sesuai. \\ \\
Kualitas kode dan dokumentasi & 8 & Kode bersih, terstruktur, terdokumentasi, dan dapat direproduksi dengan jelas. & Kode berfungsi, tetapi dokumentasi atau struktur kode masih perlu ditingkatkan. & Kode sulit dibaca, tidak terdokumentasi, atau tidak dapat direproduksi. \\}

\assessmentblock{Kuis 1 Pencarian dan Agen}{Kuis}{SCPMK0706-03101}{Kuis untuk mengukur pemahaman konsep agen cerdas dan algoritma pencarian.}{Menjawab soal pilihan ganda dan/atau uraian tentang agen dan algoritma pencarian.}{Jawaban tertulis kuis.}{Pemahaman agen cerdas dan lingkungan & 3 & Konsep agen, sifat lingkungan, dan tipe agen dijelaskan dengan benar dan lengkap. & Sebagian besar konsep dipahami, tetapi beberapa detail masih kurang tepat. & Pemahaman konsep agen tidak memadai atau penjelasan tidak relevan. \\ \\
Ketepatan algoritma pencarian & 2 & BFS, DFS, dan A* dijelaskan dengan benar, termasuk kompleksitas dan kondisi optimalitasnya. & Algoritma cukup dipahami, tetapi penjelasan kompleksitas atau optimalitas masih kurang. & Algoritma tidak dipahami dengan benar atau penjelasan mengandung kesalahan mendasar. \\}

\assessmentblock{UTS Pencarian dan Representasi Pengetahuan}{UTS berbasis pemrograman}{SCPMK0706-03101, SCPMK1005-03102}{Evaluasi tengah semester mengukur kemampuan mengimplementasikan algoritma pencarian dan representasi pengetahuan.}{Mengimplementasikan algoritma atau menjawab soal analisis sesuai soal UTS.}{Kode sumber atau jawaban tertulis UTS.}{Implementasi algoritma pencarian & 8 & Algoritma pencarian diimplementasikan dengan benar, efisien, dan menghasilkan solusi optimal pada kasus yang diberikan. & Implementasi berjalan, tetapi efisiensi atau optimasi masih belum maksimal. & Implementasi tidak menghasilkan solusi yang benar atau tidak sesuai algoritma. \\ \\
Representasi pengetahuan dan inferensi & 7 & Representasi pengetahuan (logika/aturan) dirancang dengan benar dan mesin inferensi berjalan sesuai ekspektasi. & Representasi pengetahuan ada, tetapi inferensi masih kurang tepat atau tidak konsisten. & Representasi pengetahuan tidak benar atau mesin inferensi tidak berfungsi. \\}

\assessmentblock{Kuis 2 Sistem Pakar dan ML}{Kuis}{SCPMK1005-03102}{Kuis untuk mengukur pemahaman sistem pakar, machine learning, dan metrik evaluasi model.}{Menjawab soal tentang sistem pakar, ML supervised, dan evaluasi model.}{Jawaban tertulis kuis.}{Pemahaman sistem pakar dan ML & 3 & Sistem pakar, mesin inferensi berbasis aturan, dan konsep ML supervised dijelaskan dengan benar. & Sebagian besar konsep dipahami, tetapi beberapa aspek masih kurang tepat. & Konsep sistem pakar atau ML tidak dipahami dengan benar. \\ \\
Ketepatan metrik evaluasi model & 2 & Metrik evaluasi (akurasi, presisi, recall, F1, dll.) dijelaskan dengan benar dan penerapannya pada kasus tepat. & Metrik evaluasi cukup dipahami, tetapi penerapan pada kasus masih ada kekurangan. & Metrik evaluasi tidak dipahami atau diterapkan secara keliru. \\}

\assessmentblock{PBL Implementasi Sistem AI}{Project Based Learning}{SCPMK0706-03101, SCPMK1005-03102}{Mahasiswa mengimplementasikan dan mengevaluasi sistem AI terintegrasi untuk menyelesaikan permasalahan nyata.}{Mendefinisikan masalah, memilih dan mengimplementasikan algoritma AI, mengevaluasi kinerja, dan mempresentasikan.}{Sistem AI yang berjalan, kode sumber, laporan evaluasi, dan/atau bahan presentasi.}{Ketepatan solusi dan implementasi sistem AI & 9 & Sistem AI mengimplementasikan algoritma yang tepat, menghasilkan output yang benar, dan terintegrasi dengan baik. & Sistem berjalan, tetapi pemilihan algoritma atau integrasi komponen masih kurang optimal. & Sistem tidak berjalan dengan benar atau pemilihan algoritma tidak sesuai masalah. \\ \\
Evaluasi kinerja sistem & 8 & Kinerja dievaluasi dengan metrik yang tepat, analisis mendalam, dan hasil evaluasi digunakan untuk perbaikan. & Evaluasi dilakukan dengan metrik yang cukup, tetapi analisis atau tindak lanjut masih terbatas. & Evaluasi tidak dilakukan atau metrik yang digunakan tidak sesuai konteks masalah. \\ \\
Kualitas presentasi dan dokumentasi & 5 & Presentasi jelas, dokumentasi lengkap, dan keputusan desain sistem dapat dipertahankan dengan alasan yang kuat. & Presentasi cukup jelas, tetapi dokumentasi atau argumentasi desain masih perlu ditingkatkan. & Presentasi tidak dapat menjelaskan sistem atau dokumentasi sangat minim. \\}

\assessmentblock{UAS Presentasi dan Evaluasi Sistem AI}{UAS presentasi dan evaluasi}{SCPMK1005-03102}{Mahasiswa mempresentasikan dan mempertahankan sistem AI yang dikembangkan serta mengevaluasi hasilnya secara komprehensif.}{Menyiapkan presentasi, mempresentasikan sistem, menjawab pertanyaan, dan menyerahkan laporan evaluasi final.}{Bahan presentasi, sistem AI final, laporan evaluasi komprehensif.}{Kedalaman evaluasi dan analisis hasil & 10 & Evaluasi komprehensif dengan metrik yang tepat, analisis mendalam terhadap kekuatan dan kelemahan sistem, serta rekomendasi perbaikan yang konkret. & Evaluasi cukup mendalam, tetapi beberapa aspek analisis atau rekomendasi masih perlu diperkuat. & Evaluasi dangkal, metrik tidak tepat, atau tidak ada analisis bermakna terhadap hasil sistem. \\ \\
Kemampuan presentasi dan defense & 8 & Presentasi terstruktur, penjelasan teknis jelas, dan pertanyaan dijawab dengan pemahaman mendalam terhadap sistem. & Presentasi cukup terstruktur, tetapi beberapa penjelasan teknis atau jawaban kurang memadai. & Presentasi tidak terstruktur atau tidak mampu menjelaskan sistem secara teknis. \\ \\
Validitas perbandingan dan kesimpulan & 5 & Perbandingan antarsistem/algoritma valid, kesimpulan didukung data, dan kontribusi solusi dijelaskan dengan jelas. & Perbandingan cukup valid, tetapi kesimpulan atau kontribusi solusi masih kurang spesifik. & Perbandingan tidak valid atau kesimpulan tidak didukung data yang memadai. \\}

\begin{tblr}{width=\textwidth,colspec={|Q[l,m,3.2cm]|X[l,m]|},hlines,vlines,hspan=minimal,row{1-3}={bg=headgray},rowsep=1.5pt,colsep=3pt}
Jadwal Pelaksanaan: & Minggu 1--17 sesuai RPS. Setiap evaluasi memiliki RTM dan rubrik multi-kriteria. \\
Lain-Lain Yang Diperlukan: & LMS, repositori/kode sumber, perangkat lunak sesuai mata kuliah, dan perangkat presentasi. \\
Pustaka: & Mengacu pustaka utama dan pendukung pada bagian RPS. \\
\end{tblr}
